\subsection{三高分母の完全マッチング定理：今後は仮定でなく入力として使う}
\label{nm18:matching}

整数行和・pair products を満たし，$q_0=1,q_1,q_2,q_3>1$ である局所模型を考える．
正の entries の無向 graph は，完全マッチングまたは低 ray を含む三角形と孤立頂点の和に限られる．
さらに，正の basin weights を持つことと，完全マッチング型であることは同値である．
分母上限，clutching，node，対向 $D$ 条件はいずれも不要である．
証明は算術稿の \texttt{mh18:matching} に置く．

再利用時には，低 row の三つの entries がちょうど二負一正であることを最初に課す．
正の低 entry の相手を $c$，負の二つの相手を $a,b$ とすると，
正符号 graph は $0c\sqcup ab$ であり，既存の canonical Schur 表示が全ての正の三高模型を覆う．
$\beta\le-1$ と，境界 $\beta=-1$ の完全な約数表示をそのまま使える．
ただし，全ての高 $D$ が非負となる可能性はまだ解決していない．
この係数問題の残りは $b=-\beta\ge2$，$d=-D_0\ge1$，
$\alpha\ge k\ge b+1$ の非特異領域である．
特異 Schur 系は全高 $D$ 非負の候補から除かれたが，
一般の正整数局所模型の分類にはその実在する無限族を保持する．

共通の符号分類と道型排除の再利用可能な恒等式は次節にまとめる．
